%%%%%%%%%%%%%%%%%%%%%%%%%%%%%%%%%%%%%%%%%%%%%%%%%%%%%%%%%%%%%%%%%%%%%
%  sub/en-5-discussion.tex  —  5. Discussion and Conclusion
%%%%%%%%%%%%%%%%%%%%%%%%%%%%%%%%%%%%%%%%%%%%%%%%%%%%%%%%%%%%%%%%%%%%%
\section{Discussion and Conclusion}\label{sec:discussion}

% 핵심 결과를 연구 질문·이론 틀과 연결해 해석하고, 선행연구와의 관계(확장/반박),
% 과학교육 연구·실천에의 함의, 연구의 한계와 후속 연구 방향을 제시합니다.
% Research in Science Education 은 이론적·교육적 기여를 분명히 드러내는 논의를
% 중시합니다.

Interpret the findings in relation to the research questions and framework,
connect to prior work, state implications for science-education research and
practice, and acknowledge limitations and directions for future work.
